\documentclass[12pt]{beamer}
\usetheme{Pittsburgh}
\usepackage[utf8]{inputenc}
\usepackage{amsmath}
\usepackage{amsfonts}
\usepackage{amssymb}
\usepackage{graphicx}
\usepackage{listings}
\usepackage{lstautogobble} % Provides autogobble, which is useful to remove indentation based on first line
\lstset{ %
  language=Python,
  numbers=left,
  numberstyle=\tiny,
  stepnumber=1,
  numbersep=5pt,
  breaklines=true,
  autogobble=true, % Removes indentation based on first line
}
\author{Divyesh Chauhan (AU25L20002)}
\title{Half-wave plate simulation}
\date{$2^{nd}$ March, 2026}  
\begin{document}

\begin{frame}
\titlepage
\end{frame}

\begin{frame}
\tableofcontents
\end{frame}

\begin{frame}{Introduction}
\section{Introduction}
\begin{enumerate}
	\item A half wave plate is an optical element made from a bifringent material which for a particular wavelength changes the phase difference between the $x$ and $y$ components of the electric field of light by $\pi$.
	\item As a result it preserves the intensity of that particular wavelength.
	\item Due to the change in phase difference between $x$ and $y$ components by $\pi$, left circularly polarised waves become right circularly polarised waves and vice-versa. The linearly polarised waves are flipped about their slow fast axis.
	\item Conventionally we choose the slow axis to take the negative sign from $e^ {i \pi}$ factor.
\end{enumerate}
\end{frame}

\begin{frame}[allowframebreaks, fragile]{Heatmap}
\section{Heatmap}
\begin{enumerate}
	\item The script from which this code is used to produce a video of an evolving heatmap of intensity.
	\item The function I$\_$1 gives the intensity in the region before the half-wave plate.
\begin{lstlisting}
def I_1 (z,v,lamd,t):
    v_x = v [0]* np.exp (1j*2*np.pi*(c*t-z)/lamd)
    v_y = v [1]* np.exp (1j*2*np.pi*(c*t-z)/lamd)
    return np.real (v_x)**2 + np.real (v_y)**2
\end{lstlisting}
		\item The function I$\_$2 gives the intensity in the region inside the half-wave plate.
\begin{lstlisting}
def I_2 (z,v,lamd,phi_x,phi_y,t):
    v_x = v [0]* np.exp (1j*2*np.pi*c*t/lamd-phi_x*z)
    v_y = v [1]* np.exp (1j*2*np.pi*c*t/lamd-phi_y*z)
    return np.real (v_x)**2 + np.real (v_y)**2 
\end{lstlisting}
		\item The function I$\_$3 gives the intensity in the region after the half-wave plate.
\begin{lstlisting}
def I_3 (z,v,d,lamd,phi_x,phi_y,t):
    v_x = v [0]* np.exp ( 1j*2*np.pi*(c*t-z+d)/lamd - phi_x*d )
    v_y = v [1]* np.exp ( 1j*2*np.pi*(c*t-z+d)/lamd - phi_y*d )
    return np.real (v_x)**2 + np.real (v_y)**2
\end{lstlisting}
	\item Using boolean operations the function 'intensity' produces a piecewise function to calculate the square of the real amplitude.
\begin{lstlisting}
def intensity (z,n_x,n_y,lamd,d,v,t):
    phi_x = 1j*2*np.pi*n_x/lamd
    phi_y = 1j*2*np.pi*n_y/lamd
    return (z<=0)*I_1 (z,v,lamd,t) + np.logical_and (z>0,z<d)*I_2 (z,v,lamd,phi_x,phi_y,t) + (z>=d)* I_3 (z,v,d,lamd,phi_x,phi_y,t) 
\end{lstlisting}
	\item The rotation angle doesn't impact the intensity so we will ignore it.
\end{enumerate}
\end{frame}

\begin{frame}[allowframebreaks, fragile]{Real components}
\section{Real components}
\begin{enumerate}
	\item We plot the real vector on the graph and see how it evolves. From this we can see what happens to handedness, direction of polarization etc.
	\item Let $\lambda$ be the wavelength such that with ---
	$$\lambda\ =\ \frac{d}{2\ \Delta n}$$
	\item Define $\alpha = \frac{\lambda}{\lambda^\prime}$ where $\lambda^\prime$ is the arbitrary wavelength. Then---
	$$\frac{2 \pi d}{\lambda^\prime} = \pi \alpha$$
	\item This relationship will be utilised in this fragment of code and the next one as well.
\begin{lstlisting}
def update (frame):
    v_updated = v* np.exp (1j*2*np.pi*frame/tot_frame)
    c = np.cos (theta)
    s = np.sin (theta)
    ax.clear ()
    ax.spines['left'].set_position('center')
    ax.spines['bottom'].set_position('center')
    ax.spines['right'].set_visible(False)
    ax.spines['top'].set_visible(False)
    ax.xaxis.set_ticks_position('bottom')
    ax.yaxis.set_ticks_position('left')
    if is_after == 1:
        plt.scatter (np.real (c*v_updated[0] + s*v_updated[1]), np.real ((-s*v_updated[0] + c*v_updated[1]) * np.exp (-1j*np.pi*alpha) ), color='red')
    else:
        plt.scatter (np.real (v_updated[0]), np.real (v_updated[1]) , color='red')
    ax.set_xlim (-1,1)
    ax.set_ylim (-1,1)
\end{lstlisting}	
\end{enumerate}
\end{frame}

\begin{frame}[allowframebreaks, fragile]{Complex components}
\section{Complex components}
\begin{enumerate}
	\item We plot the complex $x$ and $y$ components on the complex plane.
\begin{lstlisting}
    v_adv = v* np.exp (1j*2*np.pi*frame/tot_frame)
    v_updated = v_adv
    if is_after == 1:
        c = np.cos (theta)
        s = np.sin (theta)
        v_updated[0] = c*v_adv[0] + s*v_adv[1]
        v_updated[1] = -s*v_adv[0] + c*v_adv[1]
    ax.clear ()
    ax.spines['left'].set_position('center')
    ax.spines['bottom'].set_position('center')
    ax.spines['right'].set_visible(False)
    ax.spines['top'].set_visible(False)
    ax.xaxis.set_ticks_position('bottom')
    ax.yaxis.set_ticks_position('left')
    plt.scatter (np.real (v_updated[0]),np.imag (v_updated[0]), color='red', label= 'E_x')
    plt.scatter (np.real (v_updated[1]*np.exp (-1j*np.pi*alpha*is_after) ),np.imag (v_updated[1]*np.exp (-1j*np.pi*alpha*is_after) ), color='blue', label= 'E_y')
    plt.legend ()
    ax.set_xlim (-1,1)
    ax.set_ylim (-1,1)
\end{lstlisting}
\end{enumerate}
\end{frame}

\end{document}